\section{Spectral synthesis on the original completed tensor
source}\label{spectral-synthesis-on-the-original-completed-tensor-source}

Complete derivation CTS0--CTS6, 25 September 2026. This calculation
combines the full-source approximation GSP with the actual tensor maps
TWC. It proves statements left open when those calculations were
performed separately.

\subsection{CTS0. Corrected construction and exact
dependencies}\label{cts0.-corrected-construction-and-exact-dependencies}

The supporting datum is still
\(\tau\langle Z_1;\mathrm{no}\ Z_2\rangle\). No addition, coordinate,
metric, vector or midpoint is assigned to it. All coefficient spaces and
arithmetic operators below occur after the complete-history
reconstruction. The separate branch counters and original return
measures are unchanged. The full user argument and correction record
governs these operations; the new approximation does not supply an
alternative definition of arithmetic.

GSP7 proves continuous finite-rank operators
\(K_j=\overline K_j:\mathcal R\to\mathcal R\) converging to identity
uniformly on bounded sets in the \textbf{original} Fréchet quotient
topology. Here \[
\mathcal R=\mathcal B/\mathcal B_O=\mathcal Q/N_O,
\quad \mathcal B_O=\{F\in\mathcal B:F(\lambda)=0\ (\lambda\in\mathscr Z_O)\},
\tag{CTS0.1}
\] and \(\mathscr Z_O\) is the actual off-critical-line subset of the
nontrivial zeros of original \(\zeta\). The original full-jet quotient
\(\mathcal Q=\mathcal B/\mathcal I\) remains, with every higher jet in
the specified kernel of \(\mathcal Q\to\mathcal R\). GSP does not assume
RH, simplicity of zeros, or identification of this quotient topology
with a sequence topology.

For the original value isolator \(v_\lambda\in\mathcal R\), the formula
is \[
K_jx=\sum_{\substack{\lambda\in\mathscr Z_O\\|\Im\lambda|<T_j}}
e^{\lambda^2/j}\,\mathrm{ev}_\lambda(x)\,v_\lambda.
\tag{CTS0.2}
\] The finite sum, its nonzero Gaussian coefficients and the convergence
in the original quotient topology are proved in GSP4 and GSP7, rather
than postulated here. The heights are GSP's zero-free heights
\(T_j\in[j^6,j^6+1]\), with the finitely many initial choices as
specified there.

TWC2--3 construct the continuous map \[
b_k:\mathscr R_k=\mathcal R^{\widehat\otimes_\pi k}\longrightarrow\mathscr H_k,
\quad
(b_kx)_{\boldsymbol\rho}=\Bigl(\prod_{a=1}^k\delta_{\rho_a}\Bigr)
\varepsilon_{\boldsymbol\rho^\#}(x),
\tag{CTS0.3}
\] where
\(\varepsilon_{\boldsymbol\lambda}=\mathrm{ev}_{\lambda_1}\widehat\otimes\cdots\widehat\otimes\mathrm{ev}_{\lambda_k}\),
\(\rho^\#=1-\overline\rho\), and \[
\delta_\rho=e^{i\Im\rho\log r}(r^{\Re\rho}-r^{1-\Re\rho}),\quad r>1.
\tag{CTS0.4}
\] Every \(\delta_\rho\) on this actual off-line set is nonzero. Its
magnitude is not given a positive uniform lower bound. The receiver
retains the norms \[
P_N(y)^2=\sum_{\boldsymbol\rho\in\mathscr Z_O^k}
\Bigl(\prod_a m_{\rho_a}\Bigr)\Bigl(\prod_a(1+|\Im\rho_a|)\Bigr)^{2N}
|y_{\boldsymbol\rho}|^2.
\tag{CTS0.5}
\]

The original source transform remains \[
\Theta a(s)=\tfrac12\int_0^\infty a(u)u^s\frac{du}{u},\qquad
\Theta^{-1}F(u)=\frac{u^{-1/2}}\pi\int_{\mathbb R}F(1/2+it)u^{-it}\,dt.
\tag{CTS0.6}
\] Its full source multiplier is \[
F_0(s)=\frac{s(s-1)}8\pi^{-s/2}\Gamma(s/2)\zeta(s),
\quad F_0(0)=F_0(1)=\tfrac18,\quad F_0(-1)=F_0(2)=\tfrac\pi{24},
\] \[
F_0(-2a)=\frac{a(2a+1)(-1)^a\pi^a}{2a!}\zeta'(-2a)\quad(a\ge1).
\tag{CTS0.7}
\] At a zero \(\rho\) its full germ is the product in TWC1.5, and every
derivative uses TWC1.6's complete Leibniz sum. None of these factors or
exceptional values is changed by tensoring the already constructed maps.
The original arithmetic stays \[
\zeta(s)=1+\sum_{n\ge2}n^{-s}=\prod_p(1-p^{-s})^{-1},\qquad
-\frac{\zeta'}\zeta(s)=\sum_p\sum_{a\ge1}(\log p)p^{-as}\quad(\Re s>1).
\tag{CTS0.8}
\] This proof uses the local GSP and TWC derivations, not a new reading
of a human paper. Their pinned CC and Deligne provenance remains in
their source ledgers.

\subsection{CTS1. Equicontinuity before completed
tensoring}\label{cts1.-equicontinuity-before-completed-tensoring}

Choose increasing continuous seminorms \(p_\ell\) defining the original
Fréchet topology of \(\mathcal R\). For each fixed \(p_\ell\),
convergence of \(K_jx\) implies \(\sup_jp_\ell(K_jx)<\infty\) for every
\(x\). To obtain the needed uniform estimate, consider the closed sets
\[
E_m=\{x:\sup_jp_\ell(K_jx)\le m\},\qquad \mathcal R=\bigcup_{m\ge1}E_m.
\] The Baire property gives one set with nonempty interior. Differences
of two points in an interior translate give a neighborhood \(V\) of zero
on which \(\sup_jp_\ell(K_jx)\le2m\). Choose \(d\) and \(\epsilon>0\)
such that \(p_d(x)<\epsilon\) implies \(x\in V\). Rescale \(x\), then
let the rescaling approach the boundary, to obtain \[
\sup_jp_\ell(K_jx)\le C p_d(x)\quad(x\in\mathcal R).
\tag{CTS1.1}
\] If \(p_d(x)=0\), apply the bound to arbitrary scalar multiples of
\(x\); it forces the left side to be zero. Enlarging \(d\) and \(C\)
also bounds \(p_\ell(x)\). Thus both the identity and all \(K_j\) have a
common seminorm bound. This derives the required equicontinuity from the
actual Fréchet approximation.

On the algebraic \(k\)-fold tensor product, let \(\pi_\ell\) be the
projective seminorm obtained from \(p_\ell\) in every factor. The family
is cofinal because any finite set of source seminorms is bounded by one
of the increasing family, with constants retained. Taking the infimum
over finite tensor representations in (CTS1.1) gives \[
\pi_\ell(K_j^{\otimes k}x)\le C^k\pi_d(x).
\tag{CTS1.2}
\] These maps therefore extend continuously to the Hausdorff completion
\(\mathscr R_k\), uniformly with this estimate. Their ranges remain
finite dimensional: each is contained in the finite tensor product of
\(\mathrm{im}K_j\), which is a closed finite-dimensional subspace of the
Hausdorff completed tensor space. The coordinates
\(\varepsilon_{\boldsymbol\lambda}\) separate that finite span, so its
embedding and the stated range identification are explicit.

\subsection{CTS2. Convergence and the complete source
expansion}\label{cts2.-convergence-and-the-complete-source-expansion}

On a decomposable tensor the telescoping identity is \[
K_j^{\otimes k}(x_1\otimes\cdots\otimes x_k)-x_1\otimes\cdots\otimes x_k
=\sum_{a=1}^k K_jx_1\otimes\cdots\otimes K_jx_{a-1}
\otimes(K_jx_a-x_a)\otimes x_{a+1}\otimes\cdots\otimes x_k.
\tag{CTS2.1}
\] Each summand tends to zero in every projective seminorm: the
differing factor tends to zero and every other factor is bounded in the
required source seminorm. Hence there is convergence on finite algebraic
sums. For an arbitrary completed tensor \(x\), first choose an algebraic
tensor \(x_0\) close in \(\pi_d\) for the common bound (CTS1.2) and for
the identity. Then \[
\pi_\ell((K_j^{\widehat\otimes k}-1)x)
\le(C^k+1)\pi_d(x-x_0)+\pi_\ell((K_j^{\otimes k}-1)x_0),
\tag{CTS2.2}
\] after enlarging the constants if necessary. The second term tends to
zero, and the first can be made arbitrarily small. We have proved \[
K_j^{\widehat\otimes k}x\longrightarrow x\quad\hbox{in the original completed projective topology.}
\tag{CTS2.3}
\] Equicontinuity and finite nets also give uniform convergence on
compact subsets. We do not need or infer uniform convergence on every
bounded subset of this tensor space.

Tensoring the finite formula (CTS0.2), first on simple tensors and then
on the completion by continuity, yields the exact reconstruction \[
K_j^{\widehat\otimes k}x=
\sum_{\substack{\boldsymbol\lambda\in\mathscr Z_O^k\\|\Im\lambda_a|<T_j\ (1\le a\le k)}}
\exp\!\left(\frac1j\sum_{a=1}^k\lambda_a^2\right)
\varepsilon_{\boldsymbol\lambda}(x)
\,v_{\lambda_1}\otimes\cdots\otimes v_{\lambda_k}.
\tag{CTS2.4}
\] Every cutoff, factor and evaluation is retained. In particular the
linear span of these finite tensor blocks is dense in the original
completed source.

\subsection{CTS3. The completed tensor map is
injective}\label{cts3.-the-completed-tensor-map-is-injective}

Suppose \(b_kx=0\). Equation (CTS0.3) and the nonzero factors (CTS0.4)
give \(\varepsilon_{\boldsymbol\lambda}(x)=0\) for every actual tuple.
Formula (CTS2.4) then gives \(K_j^{\widehat\otimes k}x=0\) for every
\(j\). Passing to the actual source limit (CTS2.3) yields \(x=0\).
Consequently \[
\boxed{\ker b_k=0\quad\text{on }\mathcal R^{\widehat\otimes_\pi k}.}
\tag{CTS3.1}
\] This discharges the completed-kernel question explicitly retained in
TWC3.2 and TRC3. The older statements did not assert a nonzero kernel;
they correctly retained it until this additional convergence theorem was
available. No assumption of exactness of arbitrary completed projective
tensor products is used.

TWC3 already constructs all finite-coordinate vectors in the image. Thus
\(b_k\) is a continuous injective map with dense image into
\(\mathscr H_k\). Formula (CTS3.1) alone proves neither surjectivity nor
a topological embedding. Its stronger source-weighted receiver from
TWC10.10, and the original source topology, are retained.

\subsection{CTS4. All positive transfer forms intrinsic to the tensor
source}\label{cts4.-all-positive-transfer-forms-intrinsic-to-the-tensor-source}

The genuine product pullback and weighted transfer on \(\mathscr R_k\)
are \[
U_n=T_n^{\widehat\otimes k},\qquad V_n=(nT_{1/n})^{\widehat\otimes k},\qquad V_nU_n=n^k1.
\tag{CTS4.1}
\] Both are continuous by the already constructed coefficient action.
Consider the family \(\mathcal P_k\) of \textbf{all} jointly continuous
positive semidefinite Hermitian forms on this original completed source,
linear in the first argument, such that \[
B(U_nx,y)=B(x,V_ny)\quad(n\ge1).
\tag{CTS4.2}
\] For a finite tensor isolator
\(v_{\boldsymbol\lambda}=v_{\lambda_1}\otimes\cdots\otimes v_{\lambda_k}\),
the eigenvalue is \(n^{\sum_a\lambda_a}\). Apply (CTS4.2) to
\((v_{\boldsymbol\lambda},U_nv_{\boldsymbol\lambda})\). The left side is
\(n^{2\Re\sum_a\lambda_a}B(v_{\boldsymbol\lambda},v_{\boldsymbol\lambda})\),
and the right side is
\(n^k B(v_{\boldsymbol\lambda},v_{\boldsymbol\lambda})\). Any recovered
integer \(n>1\) therefore shows that \[
\Re\sum_a\lambda_a\ne k/2
\quad\Longrightarrow\quad B(v_{\boldsymbol\lambda},v_{\boldsymbol\lambda})=0.
\tag{CTS4.3}
\] Positivity implies Cauchy--Schwarz, proved by expanding
\(B(u+tv,u+tv)\ge0\) for arbitrary complex \(t\); a zero-norm vector
pairs to zero with every vector. Thus every noncentered finite tensor
lies in every such form's radical.

Let \[
\mathscr D_k^{\rm src}=\bigcap_{\substack{\boldsymbol\lambda\in\mathscr Z_O^k\\\Re\sum_a\lambda_a=k/2}}
\ker\varepsilon_{\boldsymbol\lambda}.
\tag{CTS4.4}
\] It is closed. If \(x\) lies there, (CTS2.4) uses only noncentered
finite tensors, so every approximation lies in every form's radical.
Continuity and (CTS2.3) give the same for \(x\).

Conversely TWC10.4 constructs the positive product-transfer form \[
B_0^{\rm src}(x,y)=\sum_{\substack{\boldsymbol\rho\in\mathscr Z_O^k\\\Re\sum_a\rho_a^\#=k/2}}
\Bigl(\prod_a m_{\rho_a}\Bigr)(b_kx)_{\boldsymbol\rho}
\overline{(b_ky)_{\boldsymbol\rho}}.
\tag{CTS4.5}
\] It is jointly continuous on the original source because \(b_k\) is
continuous. Its norm is zero precisely when all centered evaluations in
(CTS4.4) are zero, since every multiplicity is positive and every
product of \(\delta\)'s is nonzero. We conclude \[
\boxed{\bigcap_{B\in\mathcal P_k}\operatorname{rad}B
=\mathscr D_k^{\rm src}
=\overline{\operatorname{span}\{v_{\boldsymbol\lambda}:\Re\sum_a\lambda_a\ne k/2\}}^{\,\mathscr R_k}.}
\tag{CTS4.6}
\] For the last equality, the span is contained in the closed
intersection of centered kernels; its Gaussian approximants reconstruct
every vector in that intersection. This extends the common-radical
theorem to \textbf{all continuous forms on the actual tensor source},
including forms not known to factor continuously through the rapid
receiver. No assertion about the topology generated by that larger form
family is required or inferred.

\subsection{CTS5. Retained original pairing and the fixed-base
continuation}\label{cts5.-retained-original-pairing-and-the-fixed-base-continuation}

For \(k=1\), every actual off-line coordinate is noncentered, so
(CTS4.6) recovers the zero positive-form family on \(\mathcal R\). For
\(k=2\), each pair \((\lambda,\lambda^\#)\) is centered, with exact
cover eigenvalue \(n^{1+2i\Im\lambda}\). The two finite tensors \[
u=v_\lambda\otimes v_{\lambda^\#},\qquad v=v_{\lambda^\#}\otimes v_\lambda
\] survive the positive common quotient whenever the actual off-line
coordinate is present. Their original complementary Weil tensor pairing
remains \[
W_R^{\otimes2}(u\pm v,u\pm v)=\pm2m_\lambda^2,
\qquad B_0^{\rm src}(u\pm v,u\pm v)=2m_\lambda^2|\delta_\lambda|^4.
\tag{CTS5.1}
\] These are TWC10.8--12's exact original-source evaluations, not two
names for the same form. The completed-source injectivity now proves
there is no additional unseen tensor kernel explaining their difference.

The next calculation prompted by this result, using the preceding
geometric work as well, is the tensor complex \textbf{over the single
existing sphere}. Its local terms are TWC8.2--4 with their Koszul signs.
Its global angular class carries one sphere-degree factor, while the
external product in TWC9 carries a product of those factors. The
fixed-base cochain calculation is being performed on that actual
complex; no fixed \(+2\) weight is assigned to a coefficient merely to
copy Deligne's conclusion. CTS proves exact source reconstruction and
detection, not vanishing of the original specialization or its derived
connecting class.

\subsection{CTS6. Proof-use and propagation
record}\label{cts6.-proof-use-and-propagation-record}

GSP4 and GSP7 supply the additional hypothesis already proved on the
original quotient: finite-rank spectral synthesis in its Fréchet
topology. CTS1--3 perform the completed tensor argument, rather than
treating TWC's former kernel qualification as a reason to stop. CTS4
uses that reconstruction to remove the earlier restriction to forms
continuous on a selected receiver. The remaining differences of topology
and of the original Weil pairing are recorded with their exact maps and
kernels.

The full proof inputs are
\href{https://github.com/KokunoYumeto/zeta-function-research-reader/blob/main/workbenches/splitzero-tandem/continuations/20260925-full-source-tensor-and-original-divisor/gct/GAUSSIAN_SPECTRAL_SYNTHESIS.md}{Gaussian
spectral synthesis},
\href{https://github.com/KokunoYumeto/zeta-function-research-reader/blob/main/workbenches/splitzero-tandem/continuations/20260925-full-source-tensor-and-original-divisor/gct/GAUSSIAN_SPECTRAL_SYNTHESIS_INDEPENDENT.md}{its
independent derivation},
\href{https://github.com/KokunoYumeto/zeta-function-research-reader/blob/main/workbenches/splitzero-tandem/continuations/20260925-full-source-tensor-and-original-divisor/gct/ACTUAL_SPECIALIZATION_TENSOR_WEIGHT_COMPARISON.md}{the
actual tensor calculation}, and
\href{https://github.com/KokunoYumeto/zeta-function-research-reader/blob/main/workbenches/splitzero-tandem/continuations/20260925-full-source-tensor-and-original-divisor/gct/TENSOR_RECEIVER_INDEPENDENT_CHECK.md}{the
independent tensor receiver proof}. The cumulative edition retains them
in full. Earlier TWC3/TRC3 and TWC10 scopes are strengthened by CTS3 and
CTS4, respectively; their earlier text is preserved for provenance.
Nothing here claims a new complete reading of CC or Deligne, an
arithmetic purity theorem, or a resolution of RH.

\subsection{Publication source and dependency
links}\label{publication-source-and-dependency-links}

The source geometry is Alain Connes and Caterina Consani, \emph{Schemes
over F1 and zeta functions},
\href{https://arxiv.org/abs/0903.2024v3}{arXiv:0903.2024v3, §5}. The
weight-control comparison is Pierre Deligne, \emph{La conjecture de
Weil. II}, Publications mathematiques de l'IHES52 (1980),137-252,
\href{https://www.numdam.org/item/PMIHES_1980__52__137_0/}{§§3.3.11
and3.6}. Deligne material was read through the identified French
transcription and separately recorded peer source-page checks, not
Deligne-authored TeX. These sources are not asserted to contain the new
programme derivations. The
\href{https://github.com/KokunoYumeto/zeta-function-research-reader/blob/36a82ecca98addb8a98b48204e349737856c7223/workbenches/splitzero-tandem/continuations/20260924-cohomology-weight-and-character-lifts/BUILD_AND_REVIEW.md}{preceding
complete proof and reading record} records inherited source versions and
actual inspection limits. The investigator's corrected construction
remains attributed in the original text above.
